% Figura 2: validação cruzada K-fold repetida R vezes (exemplo R = K = 10).
% Somente estrutura; nenhum valor numérico.
\documentclass[border=3pt]{standalone}
% Estilo comum das figuras esquemáticas (TikZ).
% Paleta Okabe--Ito (segura para daltonismo). A cor nunca é o único código:
% todo bloco de treino/teste também traz o rótulo textual.
\usepackage[T1]{fontenc}
\usepackage[utf8]{inputenc}
\usepackage[brazilian]{babel}
\usepackage{lmodern}
\usepackage{amsmath,amssymb}
\usepackage{tikz}
\usetikzlibrary{positioning,calc,decorations.pathreplacing,arrows.meta,fit,backgrounds}

\definecolor{cTreino}{HTML}{0072B2}   % azul
\definecolor{cTeste}{HTML}{E69F00}    % laranja
\definecolor{cModA}{HTML}{009E73}     % verde-azulado
\definecolor{cModB}{HTML}{CC79A7}     % púrpura-avermelhado
\definecolor{cAlerta}{HTML}{D55E00}   % vermelhão
\definecolor{cTinta}{HTML}{222222}
\definecolor{cCinza}{HTML}{6B6B6B}

\tikzset{
  every picture/.style={font=\sffamily\small, line cap=round, line join=round, text=cTinta},
  treino/.style={draw=cTreino!80!black, fill=cTreino!22, line width=0.5pt},
  teste/.style={draw=cTeste!70!black, fill=cTeste!45, line width=0.5pt},
  caixa/.style={draw=cTinta, rounded corners=2pt, line width=0.6pt, align=flush center,
                inner sep=4pt, fill=white},
  caixaA/.style={caixa, draw=cModA!80!black, fill=cModA!14},
  caixaB/.style={caixa, draw=cModB!80!black, fill=cModB!18},
  caixaT/.style={caixa, draw=cTreino!80!black, fill=cTreino!12},
  caixaTe/.style={caixa, draw=cTeste!70!black, fill=cTeste!25},
  caixaS/.style={caixa, draw=cCinza, fill=black!4},
  seta/.style={-{Stealth[length=2.4mm]}, line width=0.6pt, draw=cTinta},
  nota/.style={font=\sffamily\footnotesize, text=cCinza, align=flush left},
}

\begin{document}
\begin{tikzpicture}[x=1cm,y=1cm]
  \def\cw{0.5}   % largura da célula do painel (a)
  \def\ch{0.36}  % altura da célula do painel (a)
  \def\rs{0.5}   % passo vertical do painel (a)
  \def\gx{10.3}  % origem x da grade do painel (b)
  \def\gc{0.46}  % lado da célula da grade
  % ---------------- Painel (a)
  \node[anchor=west, font=\sffamily\small\bfseries] at (-0.1,1.45) {(a) Uma repetição ($\ell=1$): $K=10$ dobras};
  \node[anchor=west, nota] at (-0.1,1.05) {$\pi_1$ divide $\mathcal{D}$ em $K$ subconjuntos disjuntos e estratificados};
  \foreach \m in {1,...,10} { \node[font=\sffamily\scriptsize, text=cCinza] at ({1.05+(\m-0.5)*\cw},0.55) {\m}; }
  \node[anchor=east, font=\sffamily\scriptsize, text=cCinza] at (1.0,0.55) {subconj.};
  \foreach \j in {1,...,10} {
    \pgfmathsetmacro{\y}{-(\j-1)*\rs}
    \node[anchor=east, font=\sffamily\footnotesize] at (1.0,\y) {$j=\j$};
    \foreach \m in {1,...,10} {
      \pgfmathsetmacro{\x}{1.05+(\m-1)*\cw}
      \ifnum\m=\j
        \draw[teste] (\x,\y-\ch/2) rectangle ++(\cw,\ch);
        \node[font=\sffamily\scriptsize] at ({\x+\cw/2},\y) {Te};
      \else
        \draw[treino] (\x,\y-\ch/2) rectangle ++(\cw,\ch);
      \fi
    }
    \draw[seta] ({1.05+10*\cw+0.08},\y) -- ++(0.4,0);
    \node[anchor=west, font=\sffamily\footnotesize] at ({1.05+10*\cw+0.5},\y) {$\Delta_{1,\j}$};
  }
  \draw[cTeste!70!black, line width=0.9pt, rounded corners=2pt]
    (6.48,0.25) rectangle (7.45,{-9*\rs-0.25});
  \node[nota, anchor=north west, text width=7.6cm, align=flush left] at (-0.1,{-9*\rs-0.45})
    {Dobra $j$: treina nos $K-1$ subconjuntos de treino, em azul ($n_{\mathrm{tr}}=\tfrac{K-1}{K}n$) e testa no subconjunto $\mathcal{F}_{1,j}$ ($n_{\mathrm{te}}=n/K$). Testes disjuntos; dois treinos compartilham $\tfrac{K-2}{K-1}=\tfrac89$ dos seus exemplos.};
  % ---------------- Painel (b)
  \node[anchor=west, font=\sffamily\small\bfseries] at (8.8,1.45) {(b) $R=10$ repetições $\times$ $K=10$ dobras};
  \node[anchor=west, nota] at (8.8,1.05) {cada linha usa uma nova permutação $\pi_\ell$; cada célula é uma diferença $\Delta_{\ell j}$};
  \foreach \j in {1,...,10} { \node[font=\sffamily\scriptsize, text=cCinza] at ({\gx+(\j-0.5)*\gc},0.55) {\j}; }
  \node[anchor=west, font=\sffamily\scriptsize, text=cCinza] at ({\gx+10*\gc+0.05},0.55) {dobra $j$};
  \foreach \r in {1,...,10} {
    \pgfmathsetmacro{\y}{-(\r-1)*\rs}
    \node[anchor=east, font=\sffamily\footnotesize] at ({\gx-0.05},\y) {$\ell=\r$};
    \foreach \j in {1,...,10} {
      \draw[draw=cCinza, fill=black!6, line width=0.4pt] ({\gx+(\j-1)*\gc},\y-\gc/2+0.03) rectangle ++(\gc,\gc-0.06);
    }
  }
  % destaque da linha ell=1, que corresponde ao painel (a)
  \draw[cTeste!70!black, line width=0.9pt, rounded corners=1.5pt]
    ({\gx-0.03},{-\gc/2}) rectangle ({\gx+10*\gc+0.03},{\gc/2});
  \draw[seta, cTeste!70!black, draw=cTeste!70!black] (7.5,0) .. controls (8.3,0) and (8.9,0.0) .. ({\gx-0.88},0);
  % saída
  \node[anchor=north west, caixaS, text width=6.4cm, align=flush left] (saida) at (8.8,{-9*\rs-0.45})
    {$RK=100$ diferenças pareadas $\{\Delta_{\ell j}\}$\\
     $\Rightarrow \bar\Delta$ e $\hat\sigma^2$ $\Rightarrow$ $t$ reamostrado corrigido,\\
     variância $\big(\tfrac1{RK}+\tfrac{n_{\mathrm{te}}}{n_{\mathrm{tr}}}\big)\hat\sigma^2$, referência $t_{RK-1}$};
  % rodapé
  \node[anchor=north west, nota] at (-0.1,-7.05)
    {Por algoritmo: $RK=100$ treinamentos com $0{,}9\,n$ exemplos cada; entre repetições, testes e treinos se sobrepõem. Os dois algoritmos usam as mesmas $\pi_\ell$ e as mesmas dobras.};
  \draw[treino] (-0.05,-7.75) rectangle ++(0.42,0.26); \node[anchor=west, font=\sffamily\footnotesize] at (0.4,-7.62) {treino};
  \draw[teste]  (1.6,-7.75) rectangle ++(0.42,0.26); \node[anchor=west, font=\sffamily\footnotesize] at (2.05,-7.62) {teste (Te)};
\end{tikzpicture}
\end{document}
